\section{Dataset}
\subsection{Nguồn và quy mô dữ liệu}
Trong phiên bản hiện tại, ứng dụng sử dụng release \texttt{thu\_duc\_landmarks\_v1.0.0}, gồm 65 landmark và 178 directed road-path edges tại khu vực Thủ Đức. Thông tin vị trí và đường được xây dựng từ OSM/UTraffic; một số thuộc tính traffic/flood được bổ sung để chạy các scenario thử nghiệm.

\subsection{Các thuộc tính và giả định}
Các thuộc tính chính của edge được nhóm em sử dụng như sau:
\begin{center}
{\small
\begin{tabular}{l l l}
\toprule
Thuộc tính & Ý nghĩa & Nguồn/nhãn\\
\midrule
distance & Khoảng cách của edge & DERIVED\\
free-flow time & Thời gian khi đường thông thoáng & SOURCE\_BACKED/DERIVED\\
congestion & Mức độ traffic và traffic penalty & SIMULATED hoặc dữ liệu lịch sử\\
flood risk & Nguy cơ ngập của edge & SIMULATED/ASSUMPTION\\
direction & Hướng di chuyển của đường & SOURCE\_BACKED/DERIVED\\
closure status & Edge có được phép đi qua hay không & Scenario\\
\bottomrule
\end{tabular}
}
\end{center}

Dataset chính được dùng trong giao diện và benchmark bản đồ 65 node. Các giá trị traffic/flood trong fixture kiểm chứng được ghi rõ là \texttt{SIMULATED}; vì vậy kết quả dùng để đánh giá thuật toán, không phải hướng dẫn giao thông thực tế.
